\documentclass[11pt, aspectratio = 169]{beamer}

\usetheme[progressbar=frametitle]{metropolis}
\usepackage{appendixnumberbeamer}
\usepackage{booktabs}
\usepackage{mathtools}
\usepackage{amsmath,amssymb}
\usepackage{bbm} % for \mathbbm{1}
%\usepackage{fontspec} % for emojis.
%\usepackage{fontawesome5}
\usepackage[scale=2]{ccicons}
\usepackage{pgfplots}
\usepgfplotslibrary{dateplot}
\usepackage{xspace}
\usepackage{hyperref}
\usepackage{verbatim} 
\newcommand{\themename}{\textbf{\textsc{metropolis}}\xspace}
\definecolor{links}{HTML}{1B1B81}
\hypersetup{colorlinks,linkcolor=blue,urlcolor=blue}

% listings: To highlight code:
\usepackage{listings}
\lstset{
  language=R,
  basicstyle=\ttfamily\footnotesize,
  keywordstyle=\color{blue},
  commentstyle=\color{green!50!black},
  stringstyle=\color{red},
  frame=single,
  showstringspaces=false
}
\lstset{basicstyle=\ttfamily\footnotesize, frame=single}

% To comment out large sections:
\usepackage{comment}

\title{Applied Economics}
\subtitle{Multiple Linear Regression Model II}
\author{Pello Aspuru}
\date{}
\institute{Winter 2025/26}
% \titlegraphic{\hfill\includegraphics[height=1.5cm]{logo.pdf}}
\setbeamercolor{normal text}{fg=black}
\setbeamercolor{block title}{fg=teal}
\setbeamercovered{transparent} % default faintness (~15%)

\let\svdescription\description
\let\svenddescription\enddescription
\renewenvironment{description}{%
  \svdescription%
  \let\svitem\item%
  \renewcommand\item[1][]{\svitem[\textcolor{\descriptioncolor}{##1}]}%
  }{
  \svenddescription%
  }
\newcommand\descriptioncolor{teal}

%\usepackage{Sweave}
\begin{document}
%\input{AE_w3_MLR2-concordance}

\maketitle

\begin{frame}{Plan for this Lecture}

Last lecture:
\begin{itemize}
    \item Introduced multiple linear regression model (MLR)
    \item Defined population parameters $\beta$, estimators and estimates ($\hat{\beta}$)
    \item Derived Ordinary Least Squares (OLS) estimator.
\end{itemize}    
\pause
Today lecture:
\begin{enumerate}
    \item Classical MLR Assumptions MLR.1 - MLR.4
    \item OLS Properties: Unbiasedness, Consistency and Asymptotic Normality
    \item Nonlinear Functional Forms: Polynomials, Logarithmic Transformations, Group Dummies, Interaction Terms
    \item Empirical Examples in R
\end{enumerate}
\end{frame}

\begin{comment}
    \begin{frame}{Announcements}

Homework 1 will be uploaded next Friday (Oct 25). It is due on November 4, at 6 pm.

Upload your solutions \textbf{before the deadline} via moodle (only one group member has to do so) and state the names of all group members in the \href{https://docs.google.com/spreadsheets/d/1l7be-LFx-D6pHcUKUn7duspZS-jc_b5Y40W8bhxQS0g/edit?gid=0#gid=0}{spreadsheet}.

\bigskip

Issues with R? Try swirl!

\texttt{> install.packages("swirl")} \\
\texttt{> library(swirl)} \\
\texttt{> swirl()}

Learn R interactively, in R.
\end{frame}
\end{comment}

\section{MLR Assumptions}

\begin{frame}[fragile]{Motivating Example}

Economic question: does pollution \alert{cause} a decrease of housing prices in Vienna?
 
\pause

MLR model for the \alert{population}:
$\quad Y_i = \alert{\beta_0} + \alert{\beta_1} X_{1i} + \alert{\beta_2} X_{2i} + u_{i} \quad i = 1, ... N$

\begin{itemize}
    \item Outcome variable $Y \in [0,\infty]$: housing prices
    \item Regressors: pollution concentration $X_{1i} \in (0,\infty)$, house size $X_{2i} \in (0,\infty)$.
\end{itemize}
\pause
\begin{lstlisting}[language=R]
    population <- read.csv("./AE_week3_MLR2/df_pop_house_prices.csv")
    mlr_params = lm(price ~  nox + rooms, data = population)
\end{lstlisting}

\pause
Equivalent model for observed \textcolor{teal}{sample} n:
$\quad Y_i = \textcolor{teal}{\hat{\beta}_0} + \textcolor{teal}{\hat{\beta}_1} X_{1i} + \textcolor{teal}{\hat{\beta}_2} X_{2i} + \hat{u}_{i} \quad i = 1, ... n$
\pause
\begin{lstlisting}[language=R]
    n <- 100
    df <- population[sample(nrow(population), n), ]
    ols_est = lm(price ~ nox + rooms, data = df)
\end{lstlisting}
\pause
\begin{itemize}
    \item[$\rightarrow$] \alert{$\beta$} is a \alert{number} representing the population, \textcolor{teal}{$\hat{\beta}$} is a \textcolor{teal}{random variable}.
\end{itemize}
\end{frame}

\begin{frame}[fragile]{Motivating Example}
\begin{center}
    Today: \textbf{How well} does \textcolor{teal}{random variable} \textcolor{teal}{$\hat{\beta}$} represent the \alert{number} \alert{$\beta$}?
\end{center}
\pause
$$
\text{Three MLR assumptions} \xrightarrow{\quad \text{determine} \quad} \text{feasibility and properties of } \textcolor{teal}{\hat{\beta}}
$$
\pause
\begin{enumerate}
        \item Linearity in parameters $\rightarrow$ closed-form, tractable OLS estimator function.
        %\item Random sampling $\rightarrow$ Our dataset is representative of the population.
        \item No perfect Multicolinearity $\rightarrow$ to guarantee $\hat{\beta}$ is computable.
        \item Random sampling $\rightarrow$ closed form variance of $\hat{\beta}$ (next session).
        \item Exogeneity so that
        \begin{itemize}
            \item[$\rightarrow$] $\beta$ has causal interpretation. 
            \item[$\rightarrow$] $\hat{\beta}$ is a ``good'' representation of $\beta$
        \end{itemize}
\end{enumerate}
\end{frame}


\begin{comment}
    \begin{frame}[fragile]{MLR Assumptions}

Recall the difference between population parameters, estimators and estimates:

\[
\onslide<2>{
\underbrace{
\underset{\text{Parameter}}{\beta} 
  \xrightarrow{\text{determines}} \substack{\text{Population}}}_{\text{Unobserved}}}
  %
  \onslide<3>{
  \xrightarrow{\text{iid draws}} 
  %
  \underbrace{\substack{\text{Dataset}}
  \xrightarrow{\text{estimator (OLS)}} 
  \underset{\text{Estimate}}{\hat{\beta}}}_{\text{Observed}}}
\]

    \begin{lstlisting}[language=R]
    # Open population of N = 165,000 houses in vienna: 
    population <- read.csv("./AE_week3_MLR2/df_pop_house_prices.csv")
    \end{lstlisting}
\begin{lstlisting}[language=R]
    # Obtain population parameters
    mlr_params = lm(price ~ nox + rooms, data = population)
    summary(mlr_params)
    \end{lstlisting}
\end{frame}
\end{comment}

\begin{comment}
    \begin{frame}{MLR Assumptions}

Recall the difference between population parameters, estimators and estimates:

\[
\underbrace{
\underset{\text{Parameter}}{\beta} 
  \xrightarrow{\text{determines}} \substack{\text{Population}}}_{\text{Unobserved}}
  %
  \xrightarrow{\text{iid draws}} 
  %
  \underbrace{\substack{\text{Dataset}}
  \xrightarrow{\text{estimator (OLS)}} 
  \underset{\text{Estimate}}{\hat{\beta}}}_{\text{Observed}}
\]
\pause
How well does estimate $\hat{\beta}$ represent $\beta$? Depends on...
\begin{itemize}
    \item The properties of the OLS estimator.
    \item The propertiers OLS estimator, in turn, depend on \textbf{four} MLR assumptions.
\end{itemize}
\end{frame}
\end{comment}

\begin{frame}{Assumption 1 - Linear Parameters}

\begin{center}
    In \textbf{population}, regressors $\mathbf{X}_i$ determine the outcome $Y_i$ through \alert{linear} parameters $\mathbf{\beta}$:
    \begin{align*}
    \onslide<1->{Y_i &= \alert{\beta_0} + \alert{\beta_1} X_{1i} + \alert{\beta_2} X_{2i} + \cdots + \beta_k X_{ki} + u_i \quad &\text{Linear parameters: Yes}}
    \\
    \onslide<2->{Y_i &= \beta_0 + X_{1i}^{\alert{\beta_1}} + \frac{X_{2i}}{\alert{\beta_2}} + u_i \quad &\text{Non-linear parameters: No}}
    \\
    \onslide<3->{Y_i &= \beta_0 + \beta_1 \alert{ln(X_{i})} + u_i \quad &\text{Non-linear regressors: ?}}
    \end{align*}
\end{center}
%where $\beta_0, \beta_1, \ldots, \beta_k$ are the unknown regression parameters of interest and $u_i$ is an unobservable random error term.

\begin{block}{\onslide<4->{Why do we need parameters to be linear?}}
    \onslide<5->{Linear parameters enable matrix form notation of the sample model,} \onslide<6->{consequent straightforward derivation of $\hat{\beta}$ formula:}
    $$
    \onslide<5->{\textbf{Y} = \textbf{X} \times \textbf{b} + \hat{\mathbf{u}} \quad} \onslide<6->{\rightarrow \quad SSR(b) = \hat{\mathbf{u}}' \hat{\mathbf{u}}  \quad \xrightarrow{\text{\quad FOC wrt b \quad}} \quad \hat{\boldsymbol{\beta}} =
    (\mathbf{X}'\mathbf{X})^{-1}\mathbf{X}'\mathbf{Y}}
    $$
\end{block}


\end{frame}

\begin{comment}
    \begin{frame}{MLR.2: Random sampling}

\[
\underbrace{
\underset{\text{Parameter}}{\beta} 
  \xrightarrow{\text{determines}} \substack{\text{Population}}}_{\text{Unobserved}}
  %
  \alert{\xrightarrow{\text{iid draws}}}
  %
  \underbrace{\substack{\text{Dataset}}
  \xrightarrow{\text{estimator (OLS)}} 
  \underset{\text{Estimate}}{\hat{\beta}}}_{\text{Observed}}
\]

We observe an \textbf{i.i.d.} random sample of size $n$, $\{Y_i, X_{i1}, X_{i2}, ..., X_{ik}\}, i = 1, 2, ..., n$, following the population model in MLR.1.
\pause
\bigskip

 This assumption implies independence \textbf{across} observations $i$.
 $\mathbf{X}_i \perp \mathbf{X}_{i'}$ and $u_i \perp u_{i'}$ for all $i' \neq i$
\end{frame}
\end{comment}

\begin{frame}[fragile]{Assumption 2: No perfect multicollinearity}

    Multicolinearity: Regressors $X_{1}, X_{2}...$ are \textbf{highly} correlated.
    \begin{lstlisting}[language=R]
    multicolinearity = lm(nox ~ rooms, data = df)
    summary(multicolinearity)
    \end{lstlisting}
    \pause
    \alert{Perfect} Multicollinearity: a regressor $X_1$ is a \textbf{perfect} linear function of other regressors:
        $$
        X_1 = 1, \quad \quad X_1 = X_1 + \sigma \cdot X_2 \quad \quad X_1 = 1 - X_2
        $$
    \vspace{-0.6cm}
    \begin{lstlisting}[language=R]
    df$bedrooms = df$rooms - 1
    ols_est_perf_mc = lm(price ~ nox + rooms + bedrooms, data = df)
    summary(ols_est_perf_mc)
    \end{lstlisting}
\end{frame}

\begin{frame}[fragile]{Assumption 2: No perfect multicollinearity}

\begin{block}{Why do we need the no perfect multicolinearity assumption?}
    \pause
    With perfect multicolinearity, the matrix $(\mathbf{X}'\mathbf{X})$ is not invertible $\rightarrow$ OLS estimator is not computable.
\end{block}
\pause
Easy fix: delete \alert{redundant} regressors. 
\end{frame}

\begin{frame}{Assumption 3: Random Sampling}
    We observe an \textbf{i.i.d.} random sample of the population, comprised of $n$ independent observations: 
    \[
    \{Y_i, X_{i1}, X_{i2}, \ldots, X_{ik}\},\; i = 1, 2, \ldots, n
    \]
    \pause
    This assumption implies independence \alert{across} observations $i$.
    \[
    \mathbf{X}_i \perp \mathbf{X}_{i'} \text{ and } u_i \perp u_{i'} \text{ for all } i' \neq i \quad \Rightarrow \quad E(u_i u_j \mid \mathbf{X}) = 0 \quad \text{for all } i \neq j
    \]
    \begin{block}{Why do we need the Random Sampling Assumption?}
        \pause
        For the derivation of the OLS estimate $\hat{\beta}$ variance (next session).
    \end{block}
\end{frame}

\begin{frame}{Assumption 4: Exogeneity\footnote{Also called \textbf{Zero conditional expectation} assumption}}
\vspace{-0.5cm}
\begin{center}
    Conditional Mean Independence Assumption \onslide<2->{$\quad \rightarrow \quad$ \alert{Exogeneity} Assumption}
\end{center}
\vspace{-0.5cm}
\begin{align*}
        E[u_i | X_i = A+1, X_{2i}, ..., X_{ki}] = E[u_i | X_i = A, X_{2i}, ..., X_{ki}] \onslide<2->{\alert{=0}}
\end{align*}
\onslide<2->{The error term $u_i$ has conditional expected value of zero for all regressors: $$E[u_i|X_{i1}, X_{i2}, \cdots , X_{ik}] = 0$$ }
\pause
\vspace{-0.7cm}
\begin{block}{Why do we need the Exogeneity assumption?}
    \pause
    \begin{enumerate}
        \item As with conditional mean independence assumption: causal interpretation of $\beta$.
        \pause
        \item Guarantees estimate $\hat{\beta}$ is a ``good'' representation of parameter $\beta$.
    \end{enumerate}
\end{block}
\end{frame}


\section{OLS Properties}

\begin{frame}[fragile]{Recap}
$$
\text{Three MLR assumptions} \xrightarrow{\quad \text{determine} \quad} \text{feasibility and properties of } \textcolor{teal}{\hat{\beta}}
$$
\vspace{-0.5cm}
\begin{enumerate}
        \item Linearity in parameters $\rightarrow$ closed-form, tractable OLS estimator function.
        %\item Random sampling $\rightarrow$ Our dataset is representative of the population.
        \item No perfect Multicolinearity $\rightarrow$ to guarantee $\hat{\beta}$ is computable.
        \item Random Sampling $\rightarrow$ closed form variance of $\hat{\beta}$ (next session).
        \item Exogeneity so that
        \begin{itemize}
            \item[$\rightarrow$] $\beta$ has causal interpretation. 
            \pause
            \item[$\rightarrow$] $\hat{\beta}$ is a ``good'' representation of $\beta$
        \end{itemize}
\end{enumerate}
\pause
\textbf{How well} does \textcolor{teal}{random variable} \textcolor{teal}{$\hat{\beta}$} represent the \alert{number} \alert{$\beta$}? Depends on $\hat{\beta}$ properties
\vspace{-0.2cm}
    \begin{itemize}
        \item The expectation of \textcolor{teal}{$\hat{\beta}$} $\rightarrow$ Unbiasedness and consistency.
        \item The dispersion of \textcolor{teal}{$\hat{\beta}$} $\rightarrow$ Asymptotic Distribution and variance (next session).
    \end{itemize}
\end{frame}

%  With MLR.1 – MLR.4 we can derive several properties:
%- in expectation, OLS gives true value

\begin{frame}[fragile]{Unbiasedness of OLS - Theorem}
The expectation of parameter estimate $\hat{\beta}_j$ equals the true parameter $\beta_j$:
\begin{align*}
    &E(\hat{\beta}_j) = \beta_j,\;\; j = 0,1,\ldots, k 
\end{align*}
\pause
\begin{lstlisting}[language=R]
    n <- 100
    reps <- 1000
    beta_1 <- mlr_params$coefficients["nox"] # parameter
    beta_hat_vec <- numeric(reps) # Estimates
    \end{lstlisting}
\pause
Estimate $\hat{\beta}_1$ $reps = 1000$ times, each with a \textbf{different} sample of size $n = 100$.
\begin{lstlisting}[language=R]
    for (i in 1:reps) {
        df <- population[sample(nrow(population), n), ]
        fit <- lm(price ~ nox + rooms, data = df)
        beta_hat_vec[i] <- coef(fit)["nox"]}
    \end{lstlisting}    
\end{frame}

\begin{frame}[fragile]{Unbiasedness of OLS - Graphical Illustration}
    Plot all $\hat{\beta}_1$ OLS estimates, together with $\beta_1$ true value they aim to represent:
    \pause
    \begin{lstlisting}[language=R]      
    sims <- data.frame(beta_hat = beta_hat_vec)
    ggplot(sims, aes(x = beta_hat)) +
    geom_histogram(bins = 40, alpha = 0.8) +
    geom_vline(xintercept = beta_1, linetype = "dashed", linewidth = 1)
    \end{lstlisting}   
\end{frame}

\begin{frame}{Unbiasedness - Formal Derivation}
$\underset{n \times 1}{\mathbf{Y}}$ can be expressed in terms of \alert{population parameters} as:
$
\mathbf{Y} = \mathbf{X}\boldsymbol{\hat{\beta}} + \mathbf{\hat{u}} = \alert{\mathbf{X}\boldsymbol{\beta} + \mathbf{u}}
$

\pause
Decompose OLS estimator into population parameter vector $\boldsymbol{\beta}$ and a sampling error $\mathbf{\hat{u}}$:
\begin{align*}
\hat{\boldsymbol{\beta}}=(\mathbf{X}'\mathbf{X})^{-1}\mathbf{X}'\alert{\mathbf{Y}}
                    = (\mathbf{X}'\mathbf{X})^{-1}\mathbf{X}'(\alert{\mathbf{X}\boldsymbol{\beta}+\mathbf{u}}) = \boldsymbol{\beta}+(\mathbf{X}'\mathbf{X})^{-1}\mathbf{X}'\mathbf{u}
\end{align*}
\pause
Next, take the conditional expectations:
\begin{align*}
E[\hat{\boldsymbol{\beta}}|\mathbf{X}]&= E[\boldsymbol{\beta}+(\mathbf{X}'\mathbf{X})^{-1}\mathbf{X}'\mathbf{u}|\mathbf{X}]\\
                            &=  E[\boldsymbol{\beta}|\mathbf{X}] + E[(\mathbf{X}'\mathbf{X})^{-1}\mathbf{X}'\mathbf{u}|\mathbf{X}]\\
                            &=  \boldsymbol{\beta}+(\mathbf{X}'\mathbf{X})^{-1}\mathbf{X}'\underbrace{E(\mathbf{u}|\mathbf{X})}_{=\mathbf{0}_{(nx1)}} \\
                            &=\boldsymbol{\beta} \: \hspace{+1.992cm} \text{\small Exogeneity}
\end{align*}
\pause
Last, use the Law of Iterated Expectations:
$E[\hat{\boldsymbol{\beta}}]=  E[E(\hat{\boldsymbol{\beta}}|\mathbf{X})  ]                           =E[\boldsymbol{\beta}] 
                            = \boldsymbol{\beta}$
\end{frame}

\begin{frame}{Consistency of the OLS Estimator - Theorem}
The OLS estimator $\hat{\boldsymbol{\beta}}$ converges in probability towards $\boldsymbol{\beta}$:
\begin{align*}
  %&  \hat{\boldsymbol{\beta}} \xrightarrow{p} \boldsymbol{\beta} \text{ or } \\
 &   \hat{\beta_j} \xrightarrow{p} \beta_j \text{ for all } j = 0,1,...,k
\end{align*}
Proof sketch
\begin{enumerate}
    \item Take a parameter estimate random draw $\hat{\beta_j}$, estimated from a sample of size $n$.
    \item Define the gap between the estimate and the parameter: $|\hat{\beta_j} - \beta_j |$
    \item Gap is large if $\quad |\hat{\beta_j} - \beta_j |  > \epsilon$, for some positive number $\epsilon > 0$.
    \pause
    \item The probability that the gap is large is: $\quad \mathbf{P}(|\hat{\beta_j} - \beta_j |  > \epsilon) \in [0,1]$.
    \begin{itemize}
        \item Large gap not likely $\quad \rightarrow \quad \mathbf{P}(|\hat{\beta_j} - \beta_j |  > \epsilon) \approx 0$
    \end{itemize}
\end{enumerate}
\pause
As sample size goes to $\infty$, the large gap \textbf{becomes less and less likely}:
\begin{align*} \lim_{n \to \infty} \mathbf{P}(|\hat{\beta_j} - \beta_j |  > \epsilon) = 0   \text{ for all } \epsilon > 0
\end{align*}
More in appendix.
\end{frame}

\begin{frame}[fragile]{Consistency of the OLS Estimator - Graphical Illustration}
    Find $|\hat{\beta_j} - \beta_j |$ gap $reps = 1000$ times, each with a \textbf{different} sample of size $n = 100$.
    \begin{lstlisting}[language=R]      
    gaps_beta_hat_beta <- numeric(reps)
    for (i in 1:reps) {
    df <- population[sample(nrow(population), n), ]
    fit <- lm(price ~ nox + rooms, data = df)
    gap <- abs(coef(fit)["nox"]-beta_1)
    gaps_beta_hat_beta[i] <- gap}
    \end{lstlisting}   
    \pause
    Plot the gaps:
    \begin{lstlisting}[language=R]      
    sims_gaps <- data.frame(gaps = gaps_beta_hat_beta)
    ggplot(sims_gaps, aes(x = gaps)) +
    geom_histogram(bins = 40, alpha = 0.8) +
    coord_cartesian(xlim = c(0, 2000))
    \end{lstlisting}   
    \pause
    \begin{itemize}
        \item[$\rightarrow$] What happens when we set $n = 10000$?
    \end{itemize}
\end{frame}

\begin{frame}{Asymptotic normality of the OLS Estimator - Theorem}
The random variable $\hat{\boldsymbol{\beta}}$ is asymptotically normally distributed:
$$
\hat{\boldsymbol{\beta}} \approx N\left(\boldsymbol{\beta}, \frac{1}{n} \Sigma\right)
$$
\pause
Proof sketch:
\begin{enumerate}  
    \item Rewrite gap between estimator and parameter:
    $\alert{\sqrt{n}}(\hat{\boldsymbol{\beta}}-\boldsymbol{\beta}) = 
    \left(\alert{\frac{1}{n}}\mathbf{X}'\mathbf{X}\right)^{-1}\left(\alert{\frac{1}{\sqrt{n}}}\mathbf{X}'\mathbf{u}\right)$
    \pause
    \item By law of large numbers + no perfect multicolinearity 
    $\quad \frac{1}{n}\mathbf{X}'\mathbf{X} \xrightarrow{p} \underset{k \times k}{\textcolor{teal}{\mathbf{Q}}}$
    \item By central limit theorem, exogeneity: $
\left(\tfrac{1}{\sqrt{n}}\mathbf{X}'\mathbf{u}\right)
\xrightarrow{d}
\mathcal{N}\!\left(0, \textcolor{teal}{\Omega} \right)
$
    \pause
    \item Combining previous two with Slutsky theorem:
    $
    \sqrt{n}(\hat{\boldsymbol{\beta}} - \boldsymbol{\beta}) \xrightarrow[]{d} N\left(\mathbf{0}_{(k+1)}, \textcolor{teal}{\Sigma}\right)
    $
    \item Rescaling by $\sqrt{n}$, and shifting population parameter: 
    $
    \underset{k \times 1}{\hat{\boldsymbol{\beta}}}
     \approx N \left(\boldsymbol{\beta}, \frac{1}{n} \textcolor{teal}{\Sigma} \right)
    $
    \pause
    \item[$\rightarrow$] What are \textcolor{teal}{$Q$}, \textcolor{teal}{$\Omega$}, \textcolor{teal}{$\Sigma$}? \pause Next session!
\end{enumerate}

\end{frame}

\begin{frame}{Recap: OLS Assumptions and Properties}

\begin{itemize}
    \item We revisited the OLS assumptions and the properties we get with them.
\item If the assumptions hold, our estimates $\hat{\beta}$ are correct on average (unbiased) and accurate in large samples (consistent).
\pause
\item Next session:
\begin{itemize}
    \item How does the variance matrix $Var(\hat{\boldsymbol{\beta}})$ look like?
    \item Approximate the distribution of the estimator in large samples, which we will exploit for  hypotheses testing.
\end{itemize}
\end{itemize}
\end{frame}


\section{Nonlinear Functional Forms}

\begin{frame}{Nonlinear Functional Forms}

OLS Assumption 1 requires linearity only in \color{teal}\textbf{parameters}\color{black}. OLS estimation can incorporate nonlinear transformations of the \color{teal}\textbf{variables}\color{black}.

\pause
Why include nonlinear transformations?
\begin{itemize}
    \item Linearity implies constant partial effects (independent of the level of the explanatory variable). This is often implausible, e.g. in case of decreasing marginal returns.

    \item Logarithmic transformations allow to estimate (semi-)elasticities.
\end{itemize}
\end{frame}


\begin{frame}[fragile]{(i) Polynomials of a regressor}

To model non-linear effects of $X_{1i}$ on $Y_i$ include regressor \textcolor{teal}{polynomials}: $X_{1i}^2, X_{1i}^3 \dots$
\[
\quad Y_i = \hat{\beta}_0 + \hat{\beta}_1 X_{1i} + \textcolor{teal}{\hat{\beta}_2} X_{1i}^2 + \hat{\beta}_3 X_{2i} + \hat{u}_{i} \quad i = 1, ... n
\]
\pause
\begin{lstlisting}[language=R]
    df$nox_sq <- (df$nox)^2
    ols_est = lm(price ~ nox + nox_sq + rooms, data = df)
    summary(ols_est)
\end{lstlisting}

\begin{itemize}
    \item Linear model: the effect of increasing pollutant $X_{1i}$ in \alert{+1} units on $Y_i$ is \alert{always} $\beta_1$.
    \pause
    \item Quadratic model: the effect of increasing pollutant $X_{1i}$ in \alert{+1} units on $Y_i$ is $\beta_1 + 2\beta_2 X_{1i}$. \pause Depends on \alert{level} of $X_{1i}$.
    \pause
    \item[$\rightarrow$] What does $\beta_2 < 0$ represent? Increasing or decreasing costs?
\end{itemize}

\end{frame}

\begin{frame}[fragile]{(ii) Indicator Variables}
A more flexible alternative to capture \alert{nonlinear effects} is to include indicator/dummy variables for different ranges of a continuous regressor.
\pause

Example: Split the $X_{2i}$ house rooms variable into two categories: small and big.
%$Big_i = \mathbbm{1}\left\{X_{2i} > 3 \right\}$, where
\[
Is\_Big_i  =
\begin{cases}
  1 & \text{if } X_{2i} > 5 \text{ rooms},\\
  0 & \text{otherwise.}
\end{cases}
\]
\pause
\begin{lstlisting}[language=R]
    # Generate indicator variable
    df$is_big <- 1*(df$rooms>5)
\end{lstlisting}
\end{frame}

\begin{comment}
    \begin{frame}
\frametitle{(iii) Group Dummies}
We can use group dummies to capture non-linearities in the effect of experience on the wage:

We define \color{teal}\textbf{comprehensive and mutually exclusive} \color{black} categories:
\begin{align*}
    D_{1i} &= \mathbf{1}\left\{\text{Experience}_i < 2y\right\} \\
    D_{2i} &= \mathbf{1}\left\{2y \leq \text{Experience}_i < 5y\right\} \\
    D_{3i} &= \mathbf{1}\left\{5y \leq \text{Experience}_i < 10y\right\} \\
    D_{4i} &= \mathbf{1}\left\{\text{Experience}_i \geq 10y\right\} \\
\end{align*}
We have to \alert{omit one of the dummies} in the regression. Why?
\begin{align*}
Y_i = \beta_0 + \beta_1 X_i + \alpha_1 D_{1i} + \alpha_2 D_{3i} + \alpha_3 D_{4i} + u_i
\end{align*}
\end{frame}
\end{comment}

\begin{frame}[fragile]
\frametitle{(ii) Indicator Variables}
The new specification allows the effect of house size to vary non-linearly, depending on whether the house is classified as "big" or "small".
\[
Y_i = \hat{\beta}_0 + \hat{\beta}_1 X_{1i} + \alert{\hat{\beta}_2} Is\_Big_i + \hat{u}_{i} \quad i = 1, ... n
\]
\begin{lstlisting}[language=R]
    ols_est_indic = lm(price ~ nox + is_big, data = df)
    summary(ols_est_indic)
\end{lstlisting}
\pause
\begin{itemize}
    \item Previous $\hat{\beta}_2$ assumed ``adding'' more rooms to a house had a \alert{linear} effect on price.
    \item New \alert{$\hat{\beta}_2$} assumes that being a big house causes a discrete jump in its price $\quad \rightarrow \quad$ \alert{Non-linear effect!} 
\end{itemize}
\pause
\alert{$\hat{\beta}_2$} is the average difference in house price $Y_i$ between big and small houses, holding remaining regressors $X_{1i}$ (pollution) fixed.
\end{frame}

\begin{frame}[fragile]
\frametitle{(iii) Interaction Terms}
We can further enrich our model by interacting indicator $Is\_Big_i$ and linear $X_{2i}$ variables to capture differences in slope parameters across different subgroups.
\begin{align}
    Y_i &= \hat{\beta}_0 + \hat{\beta}_1 X_{1i} + \alert{\hat{\beta}_2} X_{2i} + \hat{\beta}_3 Is\_Big_i + \hat{u}_{i} \\ 
    Y_i &= \hat{\beta}_0 + \hat{\beta}_1 X_{1i} + \alert{\hat{\beta}_2} X_{2i} + \hat{\beta}_3 Is\_Big_i + \textcolor{teal}{\hat{\beta}_4} Is\_Big_i \times X_{2i} + \hat{u}_{i}
\end{align}
\pause
Model (1) assumes that size $X_{2i}$ has the same effect on all types of houses. \\
\pause
Model (2) allows the effect of size to differ across house size categories:
\begin{itemize}
    \item for small houses: \alert{$\beta_2$}, because $Is\_Big_i \times X_{2i} = 0$
    \item for big houses: $\alert{\beta_2} + \textcolor{teal}{\beta_4}$, because $Is\_Big_i \times X_{2i} = X_{2i}$
\end{itemize}
\pause
\begin{lstlisting}[language=R]
    ols_est_inter = lm(price ~ nox + rooms*is_big, data = df)
\end{lstlisting}    
\end{frame}

%\begin{frame}{(i) Polynomials of a regressor}
%Practice Problem:

%What is the partial effect of $X_i$ in a cubic model $Y_i = \beta_0 + \beta_1 X_i + \beta_2 X_i^2 + \beta_3 X_i^3 + u_i$? 
%\end{frame}

\begin{frame}[fragile]{(iv) Logarithmic Transformations}

Starting from today's motivating example, turn pollution varaible $X_{1i}  \rightarrow \textcolor{teal}{ln (X_{1i})}$
\[
\quad Y_i = \hat{\beta}_0 + \alert{\hat{\beta}_1} \textcolor{teal}{ln(X_{1i})} + \hat{\beta}_2 X_{2i} + \hat{u}_{i} \quad i = 1, ... n
\]
\pause
Differentiate the equation with respect of $X_{1i}$ to get the partial effect of $X_{1i}$ on $Y_i$:
\begin{align*}
    \frac{\partial Y_i}{\partial X_{1i}} = \alert{\hat{\beta}_1} \frac{\partial \textcolor{teal}{ln(X_{1i})}}{\partial X_{1i}} =
    \hat{\beta}_1 \frac{1}{X_{1i}}
\end{align*}
\pause
Multiply both sides by $\partial X_{1i}$:
\begin{align*}
    \partial Y_i = \alert{\hat{\beta}_1} \cdot \frac{\partial X_{1i}}{X_{1i}} 
    \quad \xrightarrow{\quad \frac{\partial X_{1i}}{X_{1i}} \quad= \quad 0.01 \quad \equiv \quad 1\% \text{ change} \quad} \quad
     \partial Y_i = \alert{\hat{\beta}_1} \cdot 0.01
\end{align*}

For a \textbf{1\%} change in the regressor $\frac{\partial X_{1i}}{X_{1i}} = 0.01$, the change in outcome is $\partial Y_i = \alert{{\hat{\beta}}_1} \cdot 0.01$.

\end{frame}

\begin{frame}[fragile]{(iv) Logarithmic Transformations}

\color{teal}\textbf{Level-Log Model:}\color{black} $Y_i=\beta_0+\beta_1ln(X_i) + u_i$
\begin{lstlisting}[language=R]
    ols_lev_log = lm(price ~ log_nox + rooms, data = df)
\end{lstlisting}
		\begin{itemize}
		\item A 1\% increase in $X_i$ is associated with a $\frac{\beta_1}{100}$ unit increase in $Y_i$ 
	\end{itemize}
\pause
\color{teal}\textbf{Log-Linear Model:} \color{black} $ln(Y_i)=\beta_0+\beta_1 X_i + u_i$
\begin{lstlisting}[language=R]
    ols_log_lev = lm(log_price ~ nox + rooms, data = df)
\end{lstlisting}
\begin{itemize}
		\item A 1 unit increase in $X_i$ is associated with a $100\beta_1$ \% increase in $Y_i$.
  		%\item This is the semi-elasticity of $Y$ with respect to $X$!
	\end{itemize}
\pause
\color{teal}\textbf{Log-Log Model:} \color{black}$ln(Y_i)=\beta_0+\beta_1ln(X_i) + u_i$
\begin{lstlisting}[language=R]
    ols_log_log = lm(log_price ~ log_nox + rooms, data = df)
\end{lstlisting}
		\begin{itemize}
		\item A 1\% increase in $X_i$ is associated with a $\beta_1$ \% increase in $Y_i$.
		\item This is the elasticity of $Y$ with respect to $X$!
	\end{itemize}
\end{frame}


\begin{frame}
\frametitle{Interpreting coefficients with logarithmic transformations}
 
\begin{center}
\begin{tabular}{|c|c|c|c|}
  % after \\: \hline or \cline{col1-col2} \cline{col3-col4} ...
  \hline
  \textbf{Model} & \textbf{Dep. Var.} & \textbf{Indep. Var.} & \textbf{Interpretation of $\beta$} \\
  \hline
  level-level & $Y_i$    & $X_i$      & $\Delta Y_i = \beta \Delta X_i$ \\
  level-log   & $Y_i$     & $\ln(X_i)$  & $\Delta Y_i = \left(\frac{\beta}{100}\right) \% \Delta X_i$\\
  log-level   & $\ln(Y_i)$ & $X_i$      & $\% \Delta Y_i = (100 \beta) \Delta X_i$ \\
  log-log     & $\ln(Y_i)$ & $\ln(X_i)$  & $\% \Delta Y_i = \beta \% \Delta X_i$ \\
  \hline
\end{tabular}
\end{center}
$\Delta X$ and $\Delta Y$ denote absolute changes (in units). \\
$\%\Delta X$ and $\%\Delta Y$ denote relative changes (in \%). 
\end{frame}

\appendix
\begin{frame}
\centering\Large APPENDIX
    
\end{frame}
\begin{frame}{Proof Sketch - Consistency of OLS}
We can decompose the OLS estimator
into the true coefficient vector and a sampling error, and write it as a function of sample means
\begin{align*}
\hat{\boldsymbol{\beta}}&=(\mathbf{X}'\mathbf{X})^{-1}\mathbf{X}'\mathbf{Y} \\
                    &= (\mathbf{X}'\mathbf{X})^{-1}\mathbf{X}'(\mathbf{X}\boldsymbol{\beta}+\mathbf{u}) \\ 
                    &= \boldsymbol{\beta}+(\mathbf{X}'\mathbf{X})^{-1}\mathbf{X}'\mathbf{u} \\
                   &= \boldsymbol{\beta}+ \left(\sum_{i=1}^n\mathbf{X_i}\mathbf{X'_i}\right)^{-1}\sum_{i=1}^n\mathbf{X_i}u_i     \\
                    &= \boldsymbol{\beta}+ \left(\frac{1}{n}\sum_{i=1}^n\mathbf{X_i}\mathbf{X'_i}\right)^{-1}\frac{1}{n}\sum_{i=1}^n\mathbf{X_i}u_i     
\end{align*}

\end{frame}
\begin{frame}{Proof Sketch - Consistency of OLS}
Next we apply rules for the probability limit, and the  Weak Law of Large Numbers.

\begin{align*}
plim(\hat{\boldsymbol{\beta}})&=                  plim(\boldsymbol{\beta}) + plim\left[\left(\frac{1}{n}\sum_{i=1}^n\mathbf{X_i}\mathbf{X'_i}\right)^{-1}\frac{1}{n}\sum_{i=1}^n\mathbf{X_i}u_i\right]   \\
& = \boldsymbol{\beta} + plim\left[\left(\frac{1}{n}\sum_{i=1}^n\mathbf{X_i}\mathbf{X'_i}\right)^{-1}\right]  
plim\left[\frac{1}{n}\sum_{i=1}^n\mathbf{X_i}u_i\right]  \\
& = \boldsymbol{\beta} + \left[ plim \left( \frac{1}{n}\sum_{i=1}^n\mathbf{X_i}\mathbf{X'_i} \right)\right]^{-1}  plim \left[\frac{1}{n}\sum_{i=1}^n\mathbf{X_i}u_i\right]  \\
& = \boldsymbol{\beta} +  \underbrace{E[\mathbf{X_i}\mathbf{X'_i}]^{-1}}_{(k+1)x(k+1)}  \underbrace{E(\mathbf{X_i}u_i)}_{\mathbf{0_{(k+1)}}}
\end{align*}
\end{frame}





\begin{frame}{Recommended Literature}
Heiß Chapter 6.1, Chapter 7
\end{frame}

\end{document}



